\section{Discussion}
\label{sec:discussion}

\paragraph{Structure of the method space.}
The Research Skill Graph reveals several structural properties of the LLM reasoning method space. Skill degree distributions follow a power law: a small number of foundational skills (chain-of-thought prompting, self-consistency, in-context learning) participate in the majority of compositions, while most skills are specialized and sparsely connected. The graph exhibits modular community structure, with clusters corresponding to prompting strategies, verification methods, search algorithms, and training-time interventions. Interestingly, the most novel designs produced by ResearchOS tend to bridge between these clusters, combining skills from different communities---an observation consistent with theories of innovation through analogical transfer \cite{hope2017accelerating}.

\paragraph{Failure cases.}
Composition fails in several characteristic ways. First, when composed skills rest on incompatible assumptions (e.g., one skill assumes access to ground-truth intermediate steps while another assumes a fully unsupervised setting), the resulting design contains internal contradictions that the validator catches in approximately 73\% of cases but misses in the remainder. Second, the system occasionally produces \emph{over-complex} designs that compose 6--8 skills into an intricate pipeline whose interaction effects are difficult to predict. These designs score well on novelty but poorly on feasibility. Third, gap-bridging sometimes introduces speculative connective mechanisms that lack empirical grounding, inflating novelty scores while reducing technical soundness---an effect visible in the $-$Validator ablation.

\paragraph{Limitations.}
Our evaluation is subject to several important limitations. \emph{Domain specificity}: we evaluate exclusively on LLM reasoning, a domain with particularly modular and composable methods. Extension to domains with less modular structure (e.g., hardware design, wet-lab biology) may require fundamentally different skill representations. \emph{Skill extraction quality}: the system's performance depends heavily on the accuracy of extracted skill schemas; our noise injection experiments show that performance degrades substantially beyond 20\% corruption. The 20\% human calibration rate, while practical, leaves room for systematic extraction errors. \emph{LLM judge limitations}: despite our multi-judge protocol, automated evaluation of research novelty remains an open problem \cite{si2024llmideas}. Our human-LLM agreement analysis shows that novelty is the dimension with the lowest inter-rater agreement, suggesting that absolute novelty scores should be interpreted with caution. \emph{Execution scale}: our small-scale execution study covers only five designs, which is insufficient to draw general conclusions about execution success rates.

\paragraph{Broader impact.}
ResearchOS could accelerate the research ideation process, potentially democratizing access to structured research planning for researchers with less experience navigating large literatures. However, we note two risks. First, widespread adoption of composition-based ideation could \emph{homogenize} research directions, as all users draw from the same skill graph and similar composition heuristics. Second, the system may create a false sense of completeness: a well-formatted Research Proposal is not a substitute for the deep domain understanding needed to execute research successfully. Generated designs should be treated as structured hypotheses to be critically evaluated, not as validated conclusions.

\paragraph{Ethics.}
We emphasize that ResearchOS outputs are \emph{candidate designs}, not finished research. Users bear responsibility for evaluating feasibility, ethical implications, and potential harms of any design before execution. The system does not assess whether a proposed method could cause harm through its application (e.g., improved reasoning applied to generating misinformation). We advocate for human-in-the-loop deployment where experts review and filter generated designs before they enter the research pipeline.
